\changecolours{ScoutGreen}
\section[Engenharia de Requisitos]{Engenharia de Requisitos de Dados: Do Texto ao Modelo}

% ------------------------------------------------------------------------------
\twocolframe[title={Técnica de Análise Gramatical do Minimundo (Heuser, 2009)},leftcol=7.0cm,rightcol=7.0cm]{%
  \alert{Como Extrair Conceitos de um Texto?}\par
  \itemseps[topsepi=0.5mm,itemsepi=1.2mm,labelsep=2mm,leftmargini=3mm]
  \begin{itemize}
    \item \textbf{Substantivos Próprios/Comuns Principais:}
      Tornam-se candidatos a \textbf{Conjuntos de Entidades} (ex: \textit{Médico}, \textit{Paciente}, \textit{Departamento}, \textit{Livro}).
    \item \textbf{Substantivos Descritivos de Propriedades:}
      Tornam-se \textbf{Atributos} (ex: \textit{CPF}, \textit{data de nascimento}, \textit{endereço residencial}, \textit{telefones}).
    \item \textbf{Verbos Transitivos de Ação/Vínculo:}
      Tornam-se \textbf{Relacionamentos} (ex: \textit{chefiar}, \textit{alocar}, \textit{matricular-se}, \textit{fornecer}).
  \end{itemize}
}{%
  \alert{Identificando Restrições Estruturais}\par
  \itemseps[topsepi=0.5mm,itemsepi=1.2mm,labelsep=2mm,leftmargini=3mm]
  \begin{itemize}
    \item \textbf{Quantificadores de Cardinalidade:}
      Expressões como \textit{``no máximo um''}, \textit{``exatamente um''} e \textit{``diversos/múltiplos''} definem $1:1$, $1:N$ e $N:M$.
    \item \textbf{Obrigatoriedade (Participação):}
      \textit{``Todo X deve...''} indica \textbf{Participação Total (Linha Dupla)}. \textit{``Um X pode opcionalmente...''} indica \textbf{Participação Parcial}.
    \item \textbf{Expressões ``É um tipo de'' (\textit{IS-A}):}
      Indicam \textbf{Especialização/Generalização (EER)}.
  \end{itemize}
}

% ------------------------------------------------------------------------------
\begin{frame}{Pipeline Metodológico de Concepção do Projeto Prático (AVP1)}
  \vspace{0.15cm}
  \begin{center}
    \tikz[scale=0.64, transform shape]{
      \node[draw=ScoutGreen, fill=ScoutGreen!15, rounded corners=4pt, thick, minimum width=2.6cm, minimum height=1.3cm, align=center, font=\bfseries\sffamily\scriptsize] (p1) at (0,0) {ETAPA 1\\Narrativa do\\Minimundo};
      \node[draw=ScoutNavy, fill=ScoutNavy!12, rounded corners=4pt, thick, minimum width=2.6cm, minimum height=1.3cm, align=center, font=\bfseries\sffamily\scriptsize] (p2) at (3.4,0) {ETAPA 2\\Entidades Fortes\\\& Atributos};
      \node[draw=ScoutGreen, fill=ScoutGreen!15, rounded corners=4pt, thick, minimum width=2.6cm, minimum height=1.3cm, align=center, font=\bfseries\sffamily\scriptsize] (p3) at (6.8,0) {ETAPA 3\\Relacionamentos\\$1:1, 1:N, N:M$};
      \node[draw=ScoutNavy, fill=ScoutNavy!12, rounded corners=4pt, thick, minimum width=2.6cm, minimum height=1.3cm, align=center, font=\bfseries\sffamily\scriptsize] (p4) at (10.2,0) {ETAPA 4\\Entidade Fraca\\\& Herança EER};
      \node[draw=amber!90!black, fill=amber!20, rounded corners=4pt, thick, minimum width=2.6cm, minimum height=1.3cm, align=center, font=\bfseries\sffamily\scriptsize] (p5) at (13.6,0) {ETAPA 5\\Esboço DER/EER\\(Peter Chen)};

      \draw[-{Stealth[length=2.5mm]}, very thick, ScoutGreen] (p1) -- (p2);
      \draw[-{Stealth[length=2.5mm]}, very thick, ScoutNavy] (p2) -- (p3);
      \draw[-{Stealth[length=2.5mm]}, very thick, ScoutGreen] (p3) -- (p4);
      \draw[-{Stealth[length=2.5mm]}, very thick, amber!90!black] (p4) -- (p5);
    }
  \end{center}
  \vspace{0.15cm}
  \begin{columns}[t]
    \begin{column}{0.48\textwidth}
      \scriptsize
      \alert{Regra de Ouro da Rastreabilidade:}\par
      Cada elemento desenhado no diagrama da \textbf{Etapa 5} deve estar respaldado por uma frase explícita na narrativa da \textbf{Etapa 1}.
    \end{column}
    \begin{column}{0.48\textwidth}
      \scriptsize
      \alert{Prevenção de Omissões:}\par
      As tabelas intermediárias (\textbf{Etapas 2, 3 e 4}) funcionam como um filtro de qualidade antes de abrir o software CASE na Aula 08.
    \end{column}
  \end{columns}
\end{frame}
